%! Author = HK
%! Date = 2026/7/2

% Preamble
% 将 article 修改为 ctexart 即可完美支持中文，建议使用 XeLaTeX 编译器编译
\documentclass[11pt]{ctexart}

% Packages
\usepackage{amsmath}
\usepackage{siunitx} % 修复 \unit{h} 找不到命令的问题
\usepackage{enumitem} % 修复 itemsep 和 topsep 找不到属性的问题


% Document
\begin{document}
\paragraph{Master Problem 的增量重参数化}

在经典 Level Bundle 方法中，外层主问题直接以对偶变量
$(\pi,\pi_0)$ 为优化变量。由于 proximal 项本质上限制的是
当前试探点相对于稳定中心（Serious Center）的偏移，因此更自然的做法是将优化变量重新参数化为增量形式。

定义统一的对偶变量
\begin{equation}
\lambda
=
\begin{bmatrix}
\pi\\
\pi_0
\end{bmatrix},
\qquad
\bar{\lambda}
=
\begin{bmatrix}
\bar{\pi}\\
\bar{\pi}_0
\end{bmatrix},
\end{equation}
其中 $\bar{\lambda}$ 表示当前稳定中心。

定义增量变量
\begin{equation}
d
=
\lambda-\bar{\lambda}
=
\begin{bmatrix}
d_{\pi}\\
d_{\pi_0}
\end{bmatrix},
\end{equation}
即
\begin{equation}
\pi=\bar{\pi}+d_{\pi},
\qquad
\pi_0=\bar{\pi}_0+d_{\pi_0}.
\end{equation}

于是，Master Problem 可重新写为
\begin{equation}
\begin{aligned}
\max_{d,L}
\quad
&
L
-
\frac{\rho}{2}\|d\|^2
\\
\text{s.t.}\quad
&
L
\le
\omega(\lambda^i)
+
(g^i)^\top
(\bar{\lambda}+d-\lambda^i),
\qquad
\forall i\in I_t,
\end{aligned}
\end{equation}
其中
\begin{equation}
g^i=
\begin{bmatrix}
g_{\pi}^i\\
g_{\pi_0}^i
\end{bmatrix}
=
\begin{bmatrix}
z^i-x\\
f^i+\theta^i-\theta
\end{bmatrix}
\end{equation}
表示第 $i$ 个 bundle cut 对应的次梯度。

进一步整理 cut 约束，可定义常数项
\begin{equation}
\beta_i
=
(g^i)^\top
(\bar{\lambda}-\lambda^i)
+
\omega(\lambda^i),
\end{equation}
于是 Master Problem 可进一步写成
\begin{equation}
\label{eq:increment_master}
\begin{aligned}
\max_{d,L}
\quad
&
L
-
\frac{\rho}{2}\|d\|^2
\\
\text{s.t.}\quad
&
L
\le
(g^i)^\top d+\beta_i,
\qquad
\forall i\in I_t.
\end{aligned}
\end{equation}

可以看到，经过变量替换之后，Master Problem 的优化变量完全变为增量
$d$，而原始对偶变量仅通过
\begin{equation}
\lambda=\bar{\lambda}+d
\end{equation}
恢复得到。由于 proximal 项仅作用于增量，因此该重参数化保持了原问题的数学等价性，同时更加符合 Bundle 方法"在当前稳定中心附近寻找局部修正"的算法思想。

\paragraph{KKT 条件}

对于增量形式的 Master Problem \eqref{eq:increment_master}，
为每一个 cut 约束引入拉格朗日乘子
\[
\mu_i \ge 0,
\qquad
i\in I_t,
\]
则对应的拉格朗日函数为
\begin{equation}
\mathcal L(d,L,\mu)
=
L
-
\frac{\rho}{2}\|d\|^2
+
\sum_{i\in I_t}
\mu_i
\left(
(g^i)^\top d
+\beta_i
-
L
\right).
\end{equation}

\paragraph{一阶最优性条件}

分别对 $d$ 与 $L$ 求偏导，可得

\begin{align}
\nabla_d \mathcal L
&=
-\rho d
+
\sum_{i\in I_t}
\mu_i g^i
=0,
\label{eq:kkt_d}
\\
\nabla_L \mathcal L
&=
1-
\sum_{i\in I_t}
\mu_i
=0.
\label{eq:kkt_L}
\end{align}

由此得到

\begin{equation}
\sum_{i\in I_t}\mu_i=1,
\end{equation}

以及

\begin{equation}
\boxed{
d^\star
=
\frac{1}{\rho}
\sum_{i\in I_t}
\mu_i g^i .
}
\label{eq:d_star}
\end{equation}

由于
\[
g^i=
\begin{bmatrix}
g_\pi^i\\
g_{\pi_0}^i
\end{bmatrix},
\]
因此增量可以进一步写成

\begin{equation}
d^\star
=
\begin{bmatrix}
d_\pi^\star\\
d_{\pi_0}^\star
\end{bmatrix}
=
\frac1\rho
\sum_{i\in I_t}
\mu_i
\begin{bmatrix}
g_\pi^i\\
g_{\pi_0}^i
\end{bmatrix},
\end{equation}

即

\begin{align}
d_\pi^\star
&=
\frac1\rho
\sum_{i\in I_t}
\mu_i g_\pi^i,
\\
d_{\pi_0}^\star
&=
\frac1\rho
\sum_{i\in I_t}
\mu_i g_{\pi_0}^i.
\end{align}

最后恢复原始对偶变量

\begin{equation}
\lambda_{t+1}
=
\bar{\lambda}
+
d^\star,
\end{equation}

即

\begin{align}
\pi_{t+1}
&=
\bar{\pi}
+
d_\pi^\star,
\\
\pi_{0,t+1}
&=
\bar{\pi}_0
+
d_{\pi_0}^\star.
\end{align}

\paragraph{互补松弛条件}

对于每一个 bundle cut，有

\begin{equation}
\mu_i
\left(
L
-
(g^i)^\top d
-
\beta_i
\right)
=
0,
\qquad
\forall i\in I_t.
\end{equation}

该条件表明，仅有在最优解处处于活跃状态（active）的 bundle cuts 才具有非零的拉格朗日乘子，其余 cuts 对最终增量方向没有贡献。


\end{document}